\section{The Header Nobody Forwards}
\label{sec:ch15-the-header-nobody-forwards}
\index{Cosmo Router!header propagation|(}%

\code{Customer.email} works. The other field, the one carrying
\code{[Authorize]}, does not, and finding out why is what this chapter is for.

Ask for a customer through the router, holding a token that the router itself
has just accepted:

\begin{minted}[fontsize=\scriptsize,breaklines]{text}
{"errors":[{"message":"Failed to fetch from Subgraph 'accounts'.","extensions":
{"errors":[{"message":"The current user is not authorized to access this resource.",
"path":["customerById"],"extensions":{"code":"AUTH_NOT_AUTHENTICATED"}}],
"serviceName":"accounts","statusCode":200}}],"data":{"customerById":null}}
\end{minted}

Now ask for it with no token at all. The answer is identical, byte for byte.
Whether the caller was authenticated made no difference whatsoever, which is
the clue.

\subsection*{Accounts never saw the token}

\index{Cosmo Router!forwards no client headers by default}%
Chapter~\ref{ch:how-a-federated-query-actually-runs} established this and read
it as a fact about tracing headers: the router forwards no client headers to a
subgraph unless told to. It is the same rule here and it costs considerably
more.

The router validated the token, built its own view of the caller, planned the
query, and then opened a fresh HTTP request to Accounts carrying none of it.
Accounts got an anonymous request, ran \code{[Authorize]} against an anonymous
principal, and refused. Correctly.

The proof is one command away. Ask Accounts directly on 5104, no router, no
token, and it gives the same error with the same code. The router's caller was
never a variable.

Two things about that first response are worth pointing at, because they are
what makes this expensive to debug. The outer message says
\code{Failed to fetch from Subgraph 'accounts'}, which names a transport
problem that did not happen. The inner one says the user is not authorised,
which is true and reads like a bug in your own rule. Neither mentions a token,
a header, or the router's configuration, and the \code{statusCode} is 200,
because as far as HTTP is concerned everything went fine.

\subsection*{One rule, and it works}

\index{headers!op: propagate}%
The fix is five lines in \code{router/config.yaml}:

\begin{minted}[fontsize=\footnotesize,breaklines]{yaml}
headers:
  all:
    request:
      - op: propagate
        named: Authorization
\end{minted}

Restart the router and the same request answers with the customer's display
name, while an anonymous one still fails exactly as before:

\begin{minted}[fontsize=\footnotesize,breaklines]{text}
{"data":{"customerById":{"displayName":"Lars Andersen"}}}
\end{minted}

That is the whole repair.

With it in place the graph has three places a request can be refused, and
figure~\ref{fig:ch15-where-a-refusal-happens} puts them in the order they
happen.

\begin{figure}[htbp]
  \centering
  % One query, and the three places along its path that can say no.
% Referenced from chapter 15. Bare tikzpicture; the figure environment,
% caption and label live at the call site.
%
% Four boxes carry the query in turn: the client that sends it, and the three
% callers that can refuse it before an answer comes back. Three of those four
% are marked. The fourth, the client, only ever receives.
\begin{tikzpicture}[
    % 13.9cm of hand-placed coordinates against a 15.5cm measure, and the widest
    % node text pushes past it. Scaled once here rather than by moving every
    % coordinate, and transform shape keeps the node borders in proportion.
    scale=0.88, transform shape,
    font=\small,
    party/.style={draw, rounded corners=2pt, minimum width=22mm,
                 minimum height=9mm, align=center, font=\scriptsize},
    checkpoint/.style={draw, rounded corners=2pt, minimum width=30mm,
                       minimum height=13mm, align=center, font=\scriptsize},
    boundary/.style={gray, dashed, rounded corners=2pt},
    % text width is set by the two longest unbreakable strings a result box
    % has to hold, not by the prose in it: an error code and a JSON fragment
    % have no break points and a narrower box simply overflows.
    result/.style={draw, rounded corners=2pt, minimum width=40mm,
                  text width=39mm, minimum height=26mm, align=center,
                  font=\tiny, fill=black!8},
    flow/.style={-{Stealth[length=2mm]}},
    lbl/.style={font=\scriptsize, align=center, inner sep=1pt},
    note/.style={font=\scriptsize, gray, align=center, inner sep=1pt},
    closing/.style={font=\scriptsize\itshape, align=center}
  ]

  % -- the four boxes the query passes through --------------------------------
  \node[party] (client) at (0,0.6) {client};

  \node[checkpoint] (door) at (3.6,0.6)
    {access\\controller\\(before planning)};
  \node[checkpoint] (field) at (8.0,0.6)
    {field\\authorizer\\(\code{fieldConfigurations},\\during resolution)};

  \draw[boundary] (1.7,-0.35) rectangle (9.9,1.65);
  \node[lbl, anchor=south] at (5.8,1.65) {\code{Cosmo Router}};

  \node[checkpoint] (auth) at (12.0,0.6) {\code{[Authorize]}\\resolver for\\\code{customerById}};

  \draw[boundary] (10.1,-0.35) rectangle (13.9,1.65);
  \node[lbl, anchor=south] at (12.0,1.65) {\code{accounts}};

  % -- the two hops -------------------------------------------------------
  \draw[flow] (client.east) -- (door.west);
  \node[lbl, above] at (1.8,0.75) {\code{Authorization}\\header: present};

  \draw[flow] (door.east) -- (field.west);

  \draw[flow] (field.east) -- (auth.west);
  \node[lbl, above] at (10.0,0.75) {\code{Authorization} header: present};
  \node[note, below, text width=32mm] at (10.0,-0.3)
    {only because of \code{headers: op: propagate}; without that rule,
     nothing on this hop, and refusal 3 fires for every caller};

  % -- the three refusal points ------------------------------------------------
  \node[result] (r1) at (3.6,-4.2)
    {\textbf{1. at the door}\\expired or invalid token\\HTTP 401\\
     \code{\{"errors":[\{}\\
     \code{"message":"unauthorized"\}]\}}\\
     no \code{data} at all};

  \node[result] (r2) at (8.0,-4.2)
    {\textbf{2. per field}\\\code{@requiresScopes} not satisfied\\HTTP 200\\
     error per field instance,\\\code{UNAUTHORIZED\_FIELD\_OR\_TYPE}\\
     denied field nulled out};

  \node[result] (r3) at (12.0,-4.2)
    {\textbf{3. inside the subgraph}\\\code{[Authorize]} refuses\\
     HTTP 200 at both hops\\client sees\\
     \code{Failed to fetch from}\\\code{Subgraph 'accounts'}};

  \draw[flow] (door.south) -- (r1.north);
  \draw[flow] (field.south) -- (r2.north);
  \draw[flow] (auth.south) -- (r3.north);

  \node[closing, text width=132mm] at (6.5,-7.1)
    {Three shapes of refusal for one query: a request rejected before it is
     planned, a single field nulled out during resolution, and a fetch that
     fails once it reaches the subgraph. Only the last one depends on whether
     the token ever leaves the router.};

\end{tikzpicture}

  \caption{Three refusals, in the order a request meets them. A token that does
  not validate is rejected before planning, with an HTTP 401 and no data. A
  field whose scopes are not satisfied is refused during resolution, with an
  HTTP 200 and the field nulled out. A rule inside a subgraph refuses after the
  fetch, and the client is told that a fetch failed. The second hop carries the
  caller's token only because of the propagation rule above; without it the
  third refusal fires for everybody.}
  \label{fig:ch15-where-a-refusal-happens}
\end{figure}

\subsection*{What that rule costs}

\index{claim forwarding!what it hands every subgraph}%
It is worth stopping on it, because the shape of that fix is easy to
mis-generalise into \enquote{turn on header propagation and move on}.

What just happened is that every subgraph in the graph now receives the
caller's identity. Any of the seven can read any claim in that token, act on
it, and answer differently because of it, and none of that is declared in a
schema, checked by a composer, or visible in a query plan. The graph gained a
channel that carries authority and that no tool in this stack models.

The narrower alternative is to propagate to the subgraphs that need it, with a
\code{subgraphs:} block instead of \code{all:}, which is one word of
configuration and considerably more honest: it makes \enquote{Accounts is
trusted with the caller's identity} a written-down fact rather than a
side-effect. Mosaic uses \code{all:} because seven services with two rules
between them would be seven copies of the same block, and because the moment a
second service wants a subgraph-side rule the narrow version has to be edited
by somebody who will not know why it was narrow. I am not confident that is the
right call. It is the one I would defend in a code review and the one I would
reverse the first time a subgraph did something clever with a claim.

The general point survives either choice. A rule enforced inside a subgraph
needs the caller to reach the subgraph, and in a federated graph the caller
does not reach anything by default. That is a good default. It also means that
half of HotChocolate's authorization support is inert in a federated
deployment until somebody notices, and what they notice first is a field that
refuses everybody with an error naming a fetch.

\subsection*{Which of the two should own a rule}

\index{authorization!where a rule belongs}%
Now that both halves work, the question the chapter opened with can be
answered, and my answer has a shape rather than a side.

Put it in the router when the rule is about the caller and not about the data:
this field is for logged-in people, this field needs the \code{read:pii} scope.
The router can then enforce it before it fetches, the rule is in the composed
schema where every team can see it, and no subgraph needs the signing key.

Put it in the subgraph when the rule needs a fact only that service has: this
customer may read their own address and not somebody else's. No directive
expresses that, because the thing being compared is in the row.

What you should not do is put a rule in the subgraph, leave it out of the
schema, and assume the router will keep the traffic away. It will not. It does
not know the rule exists, it will plan straight through it, and what your users
get is a failed fetch.

\medskip

Both halves work now, on one field each. What the composer will let past when
the rule is wrong is a different question, and it has four answers nobody
would choose.

\index{Cosmo Router!header propagation|)}%
